\documentclass[11pt]{article}

\usepackage[utf8]{inputenc}
\usepackage[T1]{fontenc}
\usepackage{lmodern}
\usepackage{geometry}
\usepackage{amsmath,amssymb,amsfonts,amsthm,mathtools}
\usepackage{bm}
\usepackage{enumitem}
\usepackage{microtype}
\usepackage{hyperref}
\usepackage{titlesec}
\usepackage{booktabs}
\usepackage{array}
\usepackage{setspace}
\usepackage{mathrsfs}
\usepackage{physics}

\geometry{margin=1in}
\hypersetup{colorlinks=true, linkcolor=black, urlcolor=blue, citecolor=black}
\setlist[enumerate]{topsep=0.4em,itemsep=0.35em}
\setlist[itemize]{topsep=0.3em,itemsep=0.25em}
\titleformat{\section}{\large\bfseries}{\thesection.}{0.5em}{}
\titleformat{\subsection}{\normalsize\bfseries}{\thesubsection.}{0.5em}{}
\setstretch{1.06}

\newcommand{\Aether}{\AE{}ther}
\newcommand{\Aflow}{\AE{}ther-flow}
\newcommand{\TheoryName}{\Aflow{} Interpretation of Relativity}
\newcommand{\TheoryProgram}{\Aether{} / \Aflow{} framework}
\newcommand{\Sub}{\mathcal{A}}
\newcommand{\BridgeMetric}{\mathfrak{g}}

\newtheorem{definition}{Definition}
\newtheorem{proposition}{Proposition}
\newtheorem{remark}{Remark}
\newtheorem{corollary}{Corollary}
\newtheorem{theorem}{Theorem}

\title{\textbf{The \TheoryName{}}\\[0.4em]
\large Higher-Derivative Quartic Branch Obstruction to Automatic Admissible-Background Extension Beyond Uniform Control}
\author{}
\date{}

\begin{document}

\maketitle

\begin{abstract}
This manuscript performs the next disciplined step after the first broader admissible-background theorem for the higher-derivative quartic branch has been obtained on uniformly controlled bounded-geometry regions. The question is no longer whether a first theorem-level statement exists. It does. The next question is narrower: does that theorem extend automatically to a wider admissible-background class once the uniform-control hypotheses are dropped?

The answer is negative at the present scope. The current proof depends on five quantitative control channels: a smooth nondegenerate observer structure, uniform positivity of the surviving auxiliary blocks, bounded geometry relative to \(M_*\), finite-overlap patching with derivative-controlled partition functions, and exclusion of the ultrasoft observer-accessible corner through a uniform recovery-compatible hierarchy. If any one of those channels fails, the present theorem loses one of its essential estimates. The rest-space operator \(D_a\) and local momentum \(k_\perp\) cease to be uniformly well defined, or the Gaussian auxiliary elimination loses its positive denominator bound, or the curvature and overlap commutators are no longer perturbative relative to the principal quartic operator, or the residual metric self-energy
\[
\Pi_\phi^{(4)}
=
m_S^2\sigma_0^2\,
\frac{\mathcal B_4}{m_S^2\Omega^2+\mathcal B_4}
\]
ceases to be uniformly subleading because \(\varepsilon_4\not\ll1\).

Accordingly, the current broader admissible-background theorem does not extend automatically beyond uniformly controlled regions. This is not a proof that no wider theorem exists. It is a disciplined obstruction result: any stronger theorem must supply genuinely new quantitative input replacing the lost control channels. Without that new input, the theorem on uniformly controlled regions is the current honest endpoint of the admissible-background route, and nonlinear stability remains a later problem.
\end{abstract}

\section{Introduction}

The higher-derivative quartic branch has now been carried through a disciplined sequence: flat-background exact-closure screening, bridge-metric matter coupling, exact-relativistic recovery compatibility, local curved-background continuation, a region-level admissible-background package, and then a first broader admissible-background theorem on uniformly controlled bounded-geometry regions. The present note begins where that theorem stops.

The new question is not whether the theorem on uniformly controlled regions was worth proving. It was. The question is whether that theorem can now simply be read as evidence for a much wider claim on arbitrary admissible backgrounds. At the present scientific standard, the answer must be no unless one can show that the proof survives loss of the control hypotheses it actually uses.

\paragraph{Framework and claim boundary.}
\TheoryProgram{} denotes the broader \Aether{} / \Aflow{} framework built on the claims that the \Aether{} is the underlying four-dimensional substrate of reality, the \Aflow{} is its intrinsic ordered motion, observed three-dimensional space is the local experiential slice of that deeper substrate, S-time is the experienced order of change arising from matter, light, and the \Aflow{}, and observed expansion is the three-dimensional appearance of deeper four-dimensional ordered motion. Within that framework, \TheoryName{} denotes the exact-closure theory in which the effective gravitational dynamics are adopted to be exactly Einsteinian with universal matter coupling. In the active sequence, \emph{adoption} means use of that established relativistic dynamics without claiming substrate derivation, while \emph{derivation} is reserved for a first-principles recovery from explicit substrate variables.

The present note stays within that boundary. It does not prove that no wider admissible-background theorem can ever exist. It proves only the narrower and presently decisive point: the current theorem cannot be extended automatically once its quantitative control channels are dropped.

\section{Automatic Extension and Its Burden}

\begin{definition}[Automatic admissible-background extension]
An \emph{automatic extension} of the current quartic admissible-background theorem is a claim that the same observer-accessible conclusion continues to hold on a wider class of admissible regions by using only the already-recorded local curved-background continuation, admissible-background package, and uniformly controlled patching strategy, without adding new quantitative hypotheses or new dynamical estimates.
\end{definition}

\begin{definition}[Current theorem control channels]
The current broader theorem on uniformly controlled regions uses the following five control channels.
\begin{enumerate}
    \item a smooth timelike observer congruence with nondegenerate rest-space projector
    \begin{equation}
    \bar h_{\mu\nu}
    =
    \bar g_{\mu\nu}
    + \frac{1}{c^2}\bar u_\mu \bar u_\nu
    \label{eq:restprojector}
    \end{equation}
    \item uniform positivity of the surviving auxiliary blocks,
    \begin{equation}
    m_S^2\ge m_{S,\min}^2>0,
    \qquad
    \widehat M_{IJ}\ge m_{Q,\min}^2\delta_{IJ}
    \label{eq:positivity}
    \end{equation}
    \item bounded geometry relative to the quartic scale,
    \begin{equation}
    \sup_{\mathcal U}\frac{\mathcal R}{M_*^2}\ll 1,
    \qquad
    \sup_{\mathcal U}\frac{|\nabla\mathcal R|}{M_*^3}\ll 1
    \label{eq:geometry}
    \end{equation}
    \item finite-overlap patching with derivative-controlled partition functions
    \begin{equation}
    |D\zeta_\alpha|\lesssim L_{\mathcal R}^{-1},
    \qquad
    |D^2\zeta_\alpha|\lesssim L_{\mathcal R}^{-2}
    \label{eq:partition}
    \end{equation}
    \item a uniform recovery-compatible hierarchy
    \begin{equation}
    \varepsilon_{4,\mathcal U}
    \equiv
    \sup_{\mathcal U}
    \frac{1}{m_S^2\Omega^2}
    \left[
    \frac{\lambda_4}{M_*^2}k_\perp^4
    + \mathcal O\!\left(\frac{\mathcal R\,k_\perp^2}{M_*^2}\right)
    + \mathcal O\!\left(\frac{\mathcal R^2}{M_*^2}\right)
    \right]
    \ll 1.
    \label{eq:epsU}
    \end{equation}
\end{enumerate}
\end{definition}

\begin{remark}
Definition 2 is not extra structure invented after the fact. It is exactly the quantitative data used by the current theorem to patch the local continuation into a region-level result.
\end{remark}

\section{Five Obstruction Channels}

\begin{definition}[Control-loss channel]
Let \(\mathcal U\) be an admissible exact-closure background region. A \emph{control-loss channel} is any failure of one of the five ingredients in Definition 2 on \(\mathcal U\), namely: loss of smooth nondegenerate observer structure, loss of uniform auxiliary positivity, loss of bounded geometry relative to \(M_*\), loss of derivative-controlled finite-overlap patching, or loss of the non-ultrasoft recovery-compatible hierarchy.
\end{definition}

\begin{proposition}[Observer-structure loss obstructs the present operator control]
If the quartic branch observer structure fails to be smooth or the rest-space projector \(\bar h_{\mu\nu}\) becomes degenerate on a sequence of points in \(\mathcal U\), then the current proof strategy loses a uniform definition of the rest-space derivative \(D_a\) and the local momentum \(k_\perp\). Consequently the present local-to-region argument cannot be extended automatically across that sequence.
\end{proposition}

\begin{proof}
The current curved-background continuation and the broader theorem both organize the reduced scalar operator around the rest-space derivative,
\begin{equation}
\mathcal B_{4,\mathrm{cov}}
=
\frac{\lambda_4}{M_*^2}(-D^2)^2
+ \mathcal O\!\left(\frac{\mathcal R\,D^2}{M_*^2}\right)
+ \mathcal O\!\left(\frac{\nabla\mathcal R\,D}{M_*^3}\right)
+ \mathcal O\!\left(\frac{\mathcal R^2}{M_*^2}\right),
\label{eq:Bcov}
\end{equation}
and define the local hierarchy through \(k_\perp^2=\bar h^{ab}k_a k_b\). If \(\bar h_{\mu\nu}\) degenerates or \(\bar u^\mu\) is not smooth enough to define \(D_a\) coherently, then neither \eqref{eq:Bcov} nor the parameter \eqref{eq:epsU} is uniformly meaningful at the needed scope. The present argument therefore loses its first geometric datum before any broader theorem can be inferred.
\end{proof}

\begin{proposition}[Loss of auxiliary positivity obstructs uniform Gaussian elimination]
If along a sequence of points in \(\mathcal U\) one has
\begin{equation}
m_S^2(x_n)\to 0
\qquad\text{or}\qquad
\lambda_{\min}\!\bigl(\widehat M_{IJ}(x_n)\bigr)\to 0,
\label{eq:positivityloss}
\end{equation}
then the current proof strategy loses the uniform positivity needed to treat the \(\sigma\)- and \(q^I\)-sectors as uniformly auxiliary. In particular, the region-wise self-energy bound of the current theorem cannot be obtained by the same elimination argument.
\end{proposition}

\begin{proof}
The present observer-accessible reduction uses the scalar self-energy
\begin{equation}
\Pi_\phi^{(4)}
=
m_S^2\sigma_0^2\,
\frac{\mathcal B_4}{m_S^2\Omega^2+\mathcal B_4},
\label{eq:Piphi}
\end{equation}
and treats the \(q^I\)-sector as auxiliary through the positive lower bound on \(\widehat M_{IJ}\). If \(m_S^2\) approaches zero, the denominator in \eqref{eq:Piphi} loses the uniform positive time-derivative block that kept the quartic correction subleading in the intended regime. If \(\widehat M_{IJ}\) loses its lower bound, the \(q^I\)-sector is no longer uniformly removable by the same positive-definite estimate. Hence the current elimination and its region-wise bound do not extend automatically across \eqref{eq:positivityloss}.
\end{proof}

\begin{proposition}[Loss of bounded geometry obstructs the perturbative operator hierarchy]
If the curvature hierarchy \eqref{eq:geometry} fails on a region or along a sequence of regions, then the present proof strategy loses the perturbative separation between the principal quartic operator and the geometric remainder. Accordingly the current theorem cannot be extended automatically by the same local continuation estimate.
\end{proposition}

\begin{proof}
The current theorem depends on the separation of scales encoded in \eqref{eq:Bcov}. When \(\mathcal R/M_*^2\) or \(|\nabla\mathcal R|/M_*^3\) is no longer small, the curvature and commutator corrections are no longer perturbative relative to the principal term \((\lambda_4/M_*^2)(-D^2)^2\). At that point the present continuation theorem no longer gives a controlled reduced operator with the same hierarchical meaning. Therefore the broader theorem cannot be inferred automatically beyond bounded geometry by the current argument alone.
\end{proof}

\begin{proposition}[Loss of finite-overlap patch control obstructs region-level gluing]
If a wider admissible region does not admit a finite-overlap cover with partition functions obeying \eqref{eq:partition}, then the current proof strategy loses the uniform commutator estimate needed to patch local operators into a region-level operator. Consequently the present theorem does not extend automatically to that region by the same gluing step.
\end{proposition}

\begin{proof}
The current theorem glued the local operators by inserting a partition of unity and commuting derivatives past the cutoff functions. The resulting overlap contribution was controlled precisely because the atlas had finite overlap and the cutoff derivatives were bounded at the same scale as the background geometry. If those hypotheses fail, the overlap commutators need not remain of the same schematic size as the allowed geometric remainder. The proof then loses the step turning local control into a region-level operator estimate.
\end{proof}

\begin{proposition}[Ultrasoft leakage obstructs the observer-accessible subleading bound]
If the claimed observer-accessible mode family enters the ultrasoft corner so that \(\varepsilon_{4,\mathcal U}\not\ll1\), then the current theorem loses the bound that made the residual quartic self-energy uniformly subleading. Therefore the present theorem cannot be extended automatically to that regime.
\end{proposition}

\begin{proof}
The current region-level self-energy estimate is
\begin{equation}
\Pi_{\phi,\mathcal U}^{(4)}
=
\mathcal O\!\bigl(m_{S,\min}^2\sigma_0^2\,\varepsilon_{4,\mathcal U}\bigr).
\label{eq:PiU}
\end{equation}
That estimate is meaningful only when \(\varepsilon_{4,\mathcal U}\ll1\), namely when the time-derivative block \(m_S^2\Omega^2\) dominates the quartic and curvature-suppressed scalar pieces. In the ultrasoft corner this domination fails, the denominator in \eqref{eq:Piphi} can no longer be expanded in the same regime, and the residual branch correction is no longer guaranteed to be observer-accessibly subleading. Hence the current theorem does not extend automatically there.
\end{proof}

\section{Main Obstruction Result}

\begin{theorem}[Obstruction to automatic extension beyond uniform control]
Let \(\mathfrak C\) be any wider class of admissible exact-closure background regions that includes regions exhibiting at least one control-loss channel of Definition 3. Then the current broader admissible-background theorem for the higher-derivative quartic branch does not extend automatically from uniformly controlled bounded-geometry regions to all of \(\mathfrak C\).

More precisely, the existing proof strategy loses at least one indispensable estimate on every such region: coherent rest-space operator control, uniform auxiliary elimination, perturbative hierarchy of the reduced quartic operator, uniform finite-overlap gluing control, or observer-accessible subleading suppression of the residual scalar self-energy. Any stronger theorem on \(\mathfrak C\) therefore requires genuinely new quantitative input replacing the failed control channel.
\end{theorem}

\begin{proof}
Propositions 1--5 show separately that each control-loss channel breaks one of the five estimates used by the current theorem. Since an automatic extension is defined to use only the existing local continuation, admissible-background package, and current gluing strategy, the loss of any one of those estimates is enough to block that extension. Therefore the current theorem cannot be promoted automatically to any wider class containing regions where one of the control channels fails. A stronger theorem would require new hypotheses or new dynamical estimates not yet present in the active sequence.
\end{proof}

\begin{corollary}[Current claim boundary of the admissible-background route]
At the present disciplined scope, the broader admissible-background route is justified on uniformly controlled bounded-geometry regions and obstructed as an automatic extension beyond that class. The current manuscripts therefore support no stronger admissible-background claim without new quantitative input.
\end{corollary}

\begin{proof}
This is a direct reformulation of Theorem 1.
\end{proof}

\begin{remark}
Theorem 1 is not a universal no-go against every future broader theorem. It is a no-go against pretending that the current theorem already extends to a wider class without further work.
\end{remark}

\section{What This Note Establishes and What It Does Not}

The present note establishes the following.
\begin{enumerate}
    \item It identifies the exact control channels on which the current broader theorem depends.
    \item It proves that loss of any one of those channels blocks automatic extension of the current theorem by the present proof strategy.
    \item It records a disciplined admissible-background extension obstruction at the current scope.
\end{enumerate}

It does not establish the following.
\begin{enumerate}
    \item It does not prove that no wider admissible-background theorem can ever be found.
    \item It does not show that the uniformly controlled theorem is the absolute final endpoint of the theorem route.
    \item It does not solve nonlinear stability.
    \item It does not fix a unique full nonlinear covariant completion of the quartic branch.
    \item It does not complete a first-principles substrate derivation of Einsteinian gravity.
\end{enumerate}

\section{Conclusion}

The admissible-background route has now been carried through two disciplined stages. First, the quartic branch was promoted from a local continuation to a broader theorem on uniformly controlled bounded-geometry regions. Second, the present note shows that this theorem does not automatically extend once the very controls used in its proof are abandoned.

That is a real result, not a retreat. It fixes the boundary of the current theorem honestly. Any future attempt to widen the admissible-background class must now state clearly which lost control channel is being replaced by new quantitative input. Until such input exists, the theorem on uniformly controlled regions is the current admissible-background endpoint, and nonlinear stability remains a later burden rather than a problem already solved by the theorem route.

\end{document}
